% TOP_PROBLEM: {"id": 157, "title": "Sum of Cubes in F_3^n", "category_id": 2, "category": {"id": 2, "name": "combinatorics", "display_name": "Combinatorics", "description": "Counting problems, graph theory, discrete structures.", "slug": "combinatorics", "order_index": 2, "created_at": "2026-07-31T15:26:25.671Z"}, "status": "open"}
% FIRST_POSTED: null
% ENTRY_KIND: statement_only
% Imported from: batch03.tex; UnsolvedMath ID 157
\documentclass[11pt]{article}
\usepackage[margin=25mm]{geometry}
\usepackage{fontspec}
\setmainfont{Latin Modern Roman}
\usepackage{amsmath,amssymb,mathtools}
\usepackage{unicode-math}
\setmathfont{Latin Modern Math}
\usepackage{xurl,hyperref}
\hypersetup{colorlinks=true,urlcolor=blue}
\setcounter{secnumdepth}{0}
\setlength{\parindent}{0pt}
\setlength{\parskip}{5pt}
\setlength{\emergencystretch}{3em}
\newcommand{\sref}[2]{\hyperlink{source:#1:#2}{[S#2]}}
\newcommand{\cataloguescope}[1]{\paragraph{Recorded target.}#1}
\begin{document}
% UnsolvedMath ID 157; source catalogue code GREEN-069
\setcounter{section}{156}
\section{Sum of Cubes in $\mathbb F_3^n$}\label{157}
{\small\sffamily UnsolvedMath ID 157 · Combinatorics}
\subsection{Definitions and mathematical statement}

\paragraph{Shared notation.}$\mathbb N=\{1,2,\ldots\}$, and $\mathbb Z,\mathbb Q,\mathbb R,\mathbb C$ denote the integers, rationals, reals and complex numbers. Any other conventions are specified locally.


Let $V=\mathbb F_3^n$, where $\mathbb F_3=\mathbb Z/3\mathbb Z$, and regard $\{0,1\}^n$ as a subset of $V$. A cube means $L(\{0,1\}^n)$ for an invertible $\mathbb F_3$-linear map $L:V\to V$; no translation is part of this definition. Equivalently, for a basis $v_1,\ldots,v_n$ the cube consists of the subset sums $\sum_j\epsilon_jv_j$, $\epsilon_j\in\{0,1\}$. For subsets $B_1,\ldots,B_m\subset V$, their sum is $B_1+\cdots+B_m=\{b_1+\cdots+b_m:b_i\in B_i\}$.
\paragraph{Statement recorded in the supplied source.}
For every integer $n\geq1$ and every choice of $100$ invertible linear maps $L_1,\ldots,L_{100}\in\operatorname{GL}_n(\mathbb F_3)$, is
\[
 \sum_{i=1}^{100}L_i(\{0,1\}^n)=\mathbb F_3^n?
\]
Equivalently, is each vector of $\mathbb F_3^n$ a $0/1$ linear combination of the multiset union of any $100$ bases? \sref{157}{2}
\subsection{Short English statement}
Does adding any one hundred linear images of the binary cube in a vector space over the three-element field cover the entire vector space?
\subsection{Sources}
\begin{description}
\item[\hypertarget{source:157:1}{S1}] ULAM.AI and the UnsolvedMath Contributors, \emph{UnsolvedMath: A Curated Collection of Open Mathematics Problems}, record 157 (GREEN-069). CC BY 4.0. \url{https://huggingface.co/datasets/ulamai/UnsolvedMath}.
\item[\hypertarget{source:157:2}{S2}] Yang Yu, \emph{The Permanent Rank of a Matrix (Part Three): Note on the Additive Basis Conjecture}, arXiv:2510.01300v2 (2026), Introduction, Conjecture 1 and Theorem 3. The manuscript gives the stronger four-basis conclusion; since every cube contains zero, that conclusion implies the recorded one-hundred-cube statement. \url{https://arxiv.org/abs/2510.01300}.
\item[\hypertarget{source:157:3}{S3}] Ben Green, \emph{100 Open Problems}, maintained author list, Problem 26 and its comments, accessed 7 October 2026. \url{https://people.maths.ox.ac.uk/greenbj/papers/open-problems.pdf}.
\end{description}
\label{end:157}
\end{document}
